\documentclass[UTF8,12pt]{ctexart}
\usepackage[a4paper,margin=24mm]{geometry}
\usepackage{amsmath,amssymb,bm,graphicx,adjustbox}
\newcommand{\fitmath}[1]{\begin{adjustbox}{max width=\linewidth}$\displaystyle #1$\end{adjustbox}}
\setlength{\parindent}{0pt}
\setlength{\parskip}{.7em}
\sloppy
\allowdisplaybreaks
\begin{document}
\section*{2000 年考研数学三 · 真题}
\subsection*{2000 年全国硕士研究生招生考试试题}

\subsection*{一、填空题（本题共 5 小题，每小题 3 分，满分 15 分）}

（1）设 \(\displaystyle z=f\left(xy,\allowbreak \dfrac{\displaystyle x}{\displaystyle y}\right)+g\left(\dfrac{\displaystyle y}{\displaystyle x}\right)\)，其中 \(\displaystyle f,\allowbreak g\) 均可微，则 \(\displaystyle \dfrac{\displaystyle \partial z}{\displaystyle \partial x}=\)\_\_\_\_\_\_。


（2）\(\displaystyle \int_1^{+\infty}\dfrac{\displaystyle \mathrm dx}{\displaystyle e^x+e^{2-x}}=\)\_\_\_\_\_\_。


（3）若 4 阶矩阵 \(\displaystyle A\) 与 \(\displaystyle B\) 相似，矩阵 \(\displaystyle A\) 的特征值为 \(\displaystyle \dfrac{\displaystyle 1}{\displaystyle 2},\allowbreak \dfrac{\displaystyle 1}{\displaystyle 3},\allowbreak \dfrac{\displaystyle 1}{\displaystyle 4},\allowbreak \dfrac{\displaystyle 1}{\displaystyle 5}\)，则行列式 \(\displaystyle |B^{-1}-E|=\)\_\_\_\_\_\_。


（4）设随机变量 \(\displaystyle X\) 的概率密度为 \(\displaystyle f(x)=\begin{cases}\dfrac{\displaystyle 1}{\displaystyle 3},\allowbreak &x\in[0,\allowbreak 1],\allowbreak \\\dfrac{\displaystyle 2}{\displaystyle 9},\allowbreak &x\in[3,\allowbreak 6],\allowbreak \\0,\allowbreak &\text{其他}\end{cases}\) 若 \(\displaystyle k\) 使得 \(\displaystyle P\{X\ge k\}=\dfrac{\displaystyle 2}{\displaystyle 3}\)，则 \(\displaystyle k\) 的取值范围是\_\_\_\_\_\_。


（5）设随机变量 \(\displaystyle X\) 在区间 \(\displaystyle [-1,\allowbreak 2]\) 上服从均匀分布，随机变量 \(\displaystyle Y=\begin{cases}1,\allowbreak &X>0,\allowbreak \\0,\allowbreak &X=0,\allowbreak \\-1,\allowbreak &X<0,\allowbreak \end{cases}\) 则方差 \(\displaystyle D(Y)=\)\_\_\_\_\_\_。


\subsection*{二、选择题（本题共 5 小题，每小题 3 分，满分 15 分）}

（1）设对任意的 \(\displaystyle x\)，总有 \(\displaystyle \varphi(x)\le f(x)\le g(x)\)，且 \(\displaystyle \lim_{x\to\infty}[g(x)-\varphi(x)]=0\)，则 \(\displaystyle \lim_{x\to\infty}f(x)\)（　）。（A）存在且等于零。（B）存在但不一定为零。（C）一定不存在。（D）不一定存在。


（2）设函数 \(\displaystyle f(x)\) 在点 \(\displaystyle x=a\) 处可导，则函数 \(\displaystyle |f(x)|\) 在点 \(\displaystyle x=a\) 处不可导的充分条件是（　）。（A）\(\displaystyle f(a)=0\) 且 \(\displaystyle f'(a)=0\)。（B）\(\displaystyle f(a)=0\) 且 \(\displaystyle f'(a)\ne0\)。（C）\(\displaystyle f(a)>0\) 且 \(\displaystyle f'(a)>0\)。（D）\(\displaystyle f(a)<0\) 且 \(\displaystyle f'(a)<0\)。


（3）设 \(\displaystyle \alpha_1,\allowbreak \alpha_2,\allowbreak \alpha_3\) 是四元非齐次线性方程组 \(\displaystyle AX=b\) 的三个解向量，且 \(\displaystyle r(A)=3\)，\(\displaystyle \alpha_1=(1,\allowbreak 2,\allowbreak 3,\allowbreak 4)^{\mathrm T}\)，\(\displaystyle \alpha_2+\alpha_3=(0,\allowbreak 1,\allowbreak 2,\allowbreak 3)^{\mathrm T}\)，\(\displaystyle C\) 表示任意常数，则线性方程组 \(\displaystyle AX=b\) 的通解为（　）。（A）\(\displaystyle \begin{pmatrix}1\\2\\3\\4\end{pmatrix}+C\begin{pmatrix}1\\1\\1\\1\end{pmatrix}\)。（B）\(\displaystyle \begin{pmatrix}1\\2\\3\\4\end{pmatrix}+C\begin{pmatrix}0\\1\\2\\3\end{pmatrix}\)。（C）\(\displaystyle \begin{pmatrix}1\\2\\3\\4\end{pmatrix}+C\begin{pmatrix}2\\3\\4\\5\end{pmatrix}\)。（D）\(\displaystyle \begin{pmatrix}1\\2\\3\\4\end{pmatrix}+C\begin{pmatrix}3\\4\\5\\6\end{pmatrix}\)。


（4）设 \(\displaystyle A\) 为 \(\displaystyle n\) 阶实矩阵，\(\displaystyle A^{\mathrm T}\) 为 \(\displaystyle A\) 的转置矩阵，则对于线性方程组（I）：\(\displaystyle AX=0\) 和（II）\(\displaystyle A^{\mathrm T}AX=0\)，必有（　）。（A）（II）的解是（I）的解，（I）的解也是（II）的解。（B）（II）的解是（I）的解，但（I）的解不是（II）的解。（C）（I）的解不是（II）的解，（II）的解也不是（I）的解。（D）（I）的解是（II）的解，但（II）的解不是（I）的解。


（5）在电炉上安装了 4 个温控器，其显示温度的误差是随机的。在使用过程中，只要有两个温控器显示的温度不低于临界温度 \(\displaystyle t_0\)，电炉就断电。以 \(\displaystyle E\) 表示事件“电炉断电”，而 \(\displaystyle T_{(1)}\le T_{(2)}\le T_{(3)}\le T_{(4)}\) 为 4 个温控器显示的按递增顺序排列的温度值，则事件 \(\displaystyle E\) 等于（　）。（A）\(\displaystyle \{T_{(1)}\ge t_0\}\)。（B）\(\displaystyle \{T_{(2)}\ge t_0\}\)。（C）\(\displaystyle \{T_{(3)}\ge t_0\}\)。（D）\(\displaystyle \{T_{(4)}\ge t_0\}\)。


\subsection*{三、（本题满分 6 分）}

求微分方程 \(\displaystyle y''-2y'-e^{2x}=0\) 满足条件 \(\displaystyle y(0)=1,\allowbreak y'(0)=1\) 的解。


\subsection*{四、（本题满分 6 分）}

计算二重积分 \(\displaystyle \iint_D\dfrac{\displaystyle \sqrt{x^2+y^2}}{\displaystyle \sqrt{4a^2-x^2-y^2}}\,\mathrm d\sigma\)，其中 \(\displaystyle D\) 是由曲线 \(\displaystyle y=-a+\sqrt{a^2-x^2}\;(a>0)\) 和直线 \(\displaystyle y=-x\) 围成的区域。


\subsection*{五、（本题满分 6 分）}

假设某企业在两个相互分割的市场上出售同一种产品，两个市场的需求函数分别是 \(\displaystyle p_1=18-2Q_1\)，\(\displaystyle p_2=12-Q_2\)，其中 \(\displaystyle p_1\) 和 \(\displaystyle p_2\) 分别表示该产品在两个市场的价格（单位：万元／吨），\(\displaystyle Q_1\) 和 \(\displaystyle Q_2\) 分别表示该产品在两个市场的销售量（即需求量，单位：吨），并且该企业生产这种产品的总成本函数是 \(\displaystyle C=2Q+5\)，其中 \(\displaystyle Q\) 表示该产品在两个市场的销售总量，即 \(\displaystyle Q=Q_1+Q_2\)。（1）如果该企业实行价格差别策略，试确定两个市场上该产品的销售量和价格，使该企业获得最大利润；（2）如果该企业实行价格无差别策略，试确定两个市场上该产品的销售量及其统一的价格，使该企业的总利润最大化；并比较两种价格策略下的总利润大小。


\subsection*{六、（本题满分 7 分）}

求函数 \(\displaystyle y=(x-1)e^{\dfrac{\displaystyle \pi}{\displaystyle 2}+\arctan x}\) 的单调区间和极值，并求该函数图形的渐近线。


\subsection*{七、（本题满分 6 分）}

设 \(\displaystyle I_n=\displaystyle \int_0^{\pi/4}\sin^n x\cos x\,\mathrm dx\)，\(\displaystyle n=0,\allowbreak 1,\allowbreak 2,\allowbreak \ldots\)，求 \(\displaystyle \sum_{n=0}^{\infty}I_n\)。


\subsection*{八、（本题满分 6 分）}

设函数 \(\displaystyle f(x)\) 在 \(\displaystyle [0,\allowbreak \pi]\) 上连续，且 \(\displaystyle \int_0^\pi f(x)\,\mathrm dx=0\)，\(\displaystyle \int_0^\pi f(x)\cos x\,\mathrm dx=0\)。


八、（续）试证明：在 \(\displaystyle (0,\allowbreak \pi)\) 内至少存在两个不同的点 \(\displaystyle \xi_1,\allowbreak \xi_2\)，使 \(\displaystyle f(\xi_1)=f(\xi_2)=0\)。


\subsection*{九、（本题满分 8 分）}

设向量组 \(\displaystyle \alpha_1=(a,\allowbreak 2,\allowbreak 10)^{\mathrm T}\)，\(\displaystyle \alpha_2=(-2,\allowbreak 1,\allowbreak 5)^{\mathrm T}\)，\(\displaystyle \alpha_3=(-1,\allowbreak 1,\allowbreak 4)^{\mathrm T}\)，\(\displaystyle \beta=(1,\allowbreak b,\allowbreak c)^{\mathrm T}\)。试问：当 \(\displaystyle a,\allowbreak b,\allowbreak c\) 满足什么条件时，（1）\(\displaystyle \beta\) 可由 \(\displaystyle \alpha_1,\allowbreak \alpha_2,\allowbreak \alpha_3\) 线性表示，且表示唯一？（2）\(\displaystyle \beta\) 不能由 \(\displaystyle \alpha_1,\allowbreak \alpha_2,\allowbreak \alpha_3\) 线性表示？（3）\(\displaystyle \beta\) 可由 \(\displaystyle \alpha_1,\allowbreak \alpha_2,\allowbreak \alpha_3\) 线性表示，但表示不唯一？并求出一般表达式。


\subsection*{十、（本题满分 9 分）}

设有 \(\displaystyle n\) 元实二次型 \(\displaystyle f(x_1,\allowbreak x_2,\allowbreak \ldots,\allowbreak x_n)=(x_1+a_1x_2)^2+(x_2+a_2x_3)^2+\cdots+(x_{n-1}+a_{n-1}x_n)^2+(x_n+a_nx_1)^2\)，其中 \(\displaystyle a_i\;(i=1,\allowbreak 2,\allowbreak \ldots,\allowbreak n)\) 为实数。试问：当 \(\displaystyle a_1,\allowbreak a_2,\allowbreak \ldots,\allowbreak a_n\) 满足何种条件时，二次型 \(\displaystyle f(x_1,\allowbreak x_2,\allowbreak \ldots,\allowbreak x_n)\) 为正定二次型。


\subsection*{十一、（本题满分 8 分）}

（超纲题）假设 \(\displaystyle 0.50,\allowbreak 1.25,\allowbreak 0.80,\allowbreak 2.00\) 是来自总体 \(\displaystyle X\) 的简单随机样本值。已知 \(\displaystyle Y=\ln X\) 服从正态分布 \(\displaystyle N(\mu,\allowbreak 1)\)。（1）求 \(\displaystyle X\) 的数学期望 \(\displaystyle E(X)\)（记 \(\displaystyle E(X)\) 为 \(\displaystyle b\)）；（2）求 \(\displaystyle \mu\) 的置信度为 0.95 的置信区间；（3）利用上述结果求 \(\displaystyle b\) 的置信度为 0.95 的置信区间。


\subsection*{十二、（本题满分 8 分）}

设 \(\displaystyle A,\allowbreak B\) 是两个随机事件，随机变量 \(\displaystyle X=\begin{cases}1,\allowbreak &A\text{ 出现},\allowbreak \\-1,\allowbreak &A\text{ 不出现},\allowbreak \end{cases}\quad Y=\begin{cases}1,\allowbreak &B\text{ 出现},\allowbreak \\-1,\allowbreak &B\text{ 不出现}\end{cases}\) 试证明随机变量 \(\displaystyle X\) 和 \(\displaystyle Y\) 不相关的充分必要条件是 \(\displaystyle A\) 与 \(\displaystyle B\) 相互独立。

\end{document}
